%=======================================================================
% Pipeline self-test -- English body (not a real report)
%
% Every number below is copied from docs/_report_0p5b/facts.md, which was read
% out of the real trainer_state.json files.
%=======================================================================

\section{Introduction}

This supplement is the whole English body of the self-test document. It exists to prove that
the English edition of a Chinese-only report can be built from a hand-written supplement, with
the same structure as the Chinese edition: the same sections in the same order, the same
tables, and the same figure.

The 0.5B model has 504{,}179{,}896 parameters (504.18M) in 20 layers of width 1408 with 22
heads, and was trained for three legs over the same 552.33M-token packed corpus, one epoch per
leg. The pipeline under test must render all of that without inventing anything.

\subsection{Configuration}

\begin{center}\small
\begin{tabular}{@{}>{\RaggedRight\arraybackslash}p{0.34\textwidth}>{\RaggedRight\arraybackslash}p{0.58\textwidth}@{}}
\toprule
\textbf{Item} & \textbf{Value} \\
\midrule
Parameters & \textbf{504{,}179{,}896} (504.18M) \\
Layers / width / heads & 20 / 1408 / 22 \\
Feed-forward and value width & 2816 / 2816 \\
Vocabulary & 20{,}000 (tied embeddings) \\
Retention mode & \code{chunkwise}, \code{chunk\_size} 128 \\
Learned per-head decay gate & enabled (\code{use\_gla}) \\
Growth layer norm & enabled, \code{growth\_old\_dim} 768 \\
Storage precision & \code{float32} in \code{config.json}; trained in bf16 \\
\bottomrule
\end{tabular}
\end{center}

\section{Training}

The three legs share one corpus and one schedule shape: every leg restarts a full cosine with a
peak of $1.1\times10^{-4}$ and 500 warm-up steps, rather than continuing from the previous leg's
learning-rate tail.

\subsection{Training Legs}

\begin{center}\small
\begin{tabular}{@{}llrrr@{}}
\toprule
\textbf{Leg} & \textbf{Hardware} & \textbf{Steps} & \textbf{Tokens per step} & \textbf{Total tokens} \\
\midrule
ep1 & RTX 3090 24\,GB & 67{,}390 & 8{,}192 & 552{,}058{,}880 \\
ep2 & RTX PRO 6000 & 16{,}848 & 32{,}768 & 552{,}140{,}800 \\
ep3 & RTX PRO 6000 & 16{,}848 & 32{,}768 & 552{,}140{,}800 \\
\bottomrule
\end{tabular}
\end{center}

Cumulatively the three legs consumed 1{,}656{,}340{,}480 tokens, or 3.285 tokens per parameter.

\subsection{Validation Loss}

\begin{center}\small
\begin{tabular}{@{}lrrr@{}}
\toprule
\textbf{Leg} & \textbf{First eval} & \textbf{Best eval} & \textbf{Best step} \\
\midrule
ep1 & 2.6358 & \textbf{2.1035} & 65{,}000 \\
ep2 & 2.1800 & \textbf{1.9865} & 16{,}848 \\
ep3 & 2.1174 & \textbf{1.9459} & 16{,}848 \\
\bottomrule
\end{tabular}
\end{center}

\subsection{The Cost of the Hot Restarts}

\begin{center}\small
\begin{tabular}{@{}lrrr@{}}
\toprule
\textbf{Hand-off} & \textbf{Previous best} & \textbf{Next first eval} & \textbf{Difference} \\
\midrule
ep1 $\rightarrow$ ep2 & 2.1035 @65{,}000 & 2.1800 @2{,}106 & $+0.0764$ \\
ep2 $\rightarrow$ ep3 & 1.9865 @16{,}848 & 2.1174 @2{,}106 & $+0.1310$ \\
\bottomrule
\end{tabular}
\end{center}

Both hand-offs made the model \emph{worse} before it recovered: the first evaluation of a new
leg sits above the previous leg's best, by 0.0764 and 0.1310 respectively, and both legs do
eventually finish below that starting point.

\section{Appendix A: Reproduction}

\begin{lstlisting}[language=Python]
# 1) the figures, from the real state files
python makefigs.py

# 2) Markdown -> LaTeX, then two XeLaTeX passes per edition
python mktex.py --report selftest
\end{lstlisting}

\backmattersection{References}

\begin{itemize}
  \item Retentive Network: A Successor to Transformer for Large Language Models. arXiv:2307.08621.
  \item Gated Linear Attention Transformers with Hardware-Efficient Training. arXiv:2312.06635.
\end{itemize}
